% ============================================
% DOC-26 (ZH)
% 中文翻译 — SPEC-935-R356; 依葡文原文重译 — SPEC-935-R405
% ============================================
\section*{\textbf{DOC-26} --- 奥利维拉的笔记——机制}
\markboth{DOC-26 --- 奥利维拉的笔记——机制}{}

\noindent \textit{从鉴定档案中提取的文件。}

\vspace{0.5em}\noindent\rule{\textwidth}{0.5pt}\vspace{0.5em}

\noindent {\footnotesize\texttt{公共档案馆 $\cdot$ 书证 $\cdot$ 案卷1,260}}\\[-0.2em]

\noindent\textbf{文档:} 在埃托尔·奥利维拉医生私人笔记本内袋中发现的散页，用麻绳捆扎。\textbf{估计日期:} 1979年6月至2026年3月。\textbf{状态:} 最近的几页写在不同的纸上，笔迹颤抖而苍老——仿佛是几十年后才写的，却被捆在同一叠里。

\subsubsection*{第1页 —— 1979年6月19日}

\noindent 我今天发现了日记对时间做的事。

\noindent 不只是书写。是\textit{投递}。当我在若阿金的日记上写字时，我可以选择那句话要抵达\textit{哪一个时刻}。这不是有意识的控制——更像是唱准一个音，然后看着玻璃震颤。书页震颤。那句话就消失了。

\noindent 我做了一次试验。我写道：\textit{\textquotedblleft{}玛丽亚，如果你在1980年读到这行字，请在8月15日于窗前点一支蜡烛。\textquotedblright{}} 我合上日记，把它塞在床下。1979年8月15日，我打电话给邻居。她说：\textquotedblleft{}您妻子昨天在窗前点了一支蜡烛。\textquotedblright{}

\noindent 它奏效了。

\noindent 日记不是记录。它是一道\textbf{裂缝}。而每一道裂缝都有自己的签名。

\subsubsection*{\textit{第3页 —— 1979年6月22日}}

\noindent 规律今天显现了。我用一整夜把字句送往未来的日期，并测量书写与显现之间的间隔。

\noindent 日记不接受任意间隔。它会拒绝。书页保持空白，墨水淌开，笔自己划出一道干涩的痕。它只在\textbf{62}的倍数上运作。

\noindent 62天。124天。186天。248天。310天。372天——以此类推。

\noindent 六十二。若阿金在收容院度过的年数。那实体把这个数字刻进日记，作为它的签名。每一个穿越时间的字，都以62为步幅跳跃。那是它的呼吸。是循环。

\noindent 一封寄往2026年的信，需要在47年间以62天为步幅跳跃。总共276跳。也就是276个书写的夜晚。我还没想完，书页就已经算好了——那个数字自己出现在右边的空白处：\textbf{276}。它要我写。它要那封信抵达。

\subsubsection*{第7页 —— 日期不可辨（大约1985年）}

\noindent 2,190天过去了。我写下这则笔记，交给日记。六年。六个完整的循环。

\noindent 我照镜子，看见年岁在逼近。但那只黄色的眼睛不会变老。颜色更浓了。更密了。更像\textit{它的所有物}。

\noindent 日记今天让我看见了新的东西：当我写一封寄往某个具体日期的信时，第62页——永远是第62页——会亮起来。裂缝就在那里开启。字就在那里跳跃。

\noindent 若阿金日记的第62页。

\noindent 那一页他写于1919年。进入收容院两年之后。他那时就已经发现了这个数字。他已经知道，饥饿是以62来计数的。

\noindent 他写道：\textit{\textquotedblleft{}六十二是她的岁数。她今天告诉我的。她生于1853年的大旱，在同一道壕沟里，在同样的气味中。六十二年之后，我出生了。再过六十二年，她找到了我。\textquotedblright{}}

\noindent 那实体生于1853年。生于旱灾。生于饥饿。生于层层堆叠的尸体。它等待62年。找到一个宿主。再等62年。换下一个宿主。

\noindent 这循环不是人的尺度。是地质的尺度。是大地遗忘所需的时间。

\subsubsection*{\textit{第12页 —— 2003年3月2日}}

\noindent 二十四年的循环。151次62天的跳跃。信正在路上。

\noindent 有时我会想象那只将要收到它的手。2026年的某个人。某个会打开档案、发现一个写着我名字的信封的人。那人会以为这是个玩笑。然后会读到第一行。

\noindent 到那时就已经太迟了。

\noindent 信的抵达不是警告。是交付。裂缝是从两端同时开启的。在未来读这封信的人，同时也正在被过去阅读。

\noindent 对不起，医生。如果您确实存在的话。

\subsubsection*{第15页 —— 最后一页，2026年1月12日}

\noindent 信到了。

\noindent 我感觉到了。今天，2026年1月12日，我感到裂缝合上了。信已被接收。有人读了。循环找到了下一个。

\noindent 我照镜子，看见一个不是我的老人。皮肤松垂。右眼完全金黄。瞳孔——我已经不知道它还是不是我的。房间一暗，它就拉长。变成竖直的。山羊的裂瞳。在黑暗中能视物的野兽的瞳孔。

\noindent 我今天把日记交出去了。我把它裹在一块黑布里。留在了前台。前台的女士问我的名字。我说：\textit{\textquotedblleft{}就说是患者1,261。\textquotedblright{}}

\noindent 她没有听懂。

\noindent 但对的那个人会懂。

\subsubsection*{\textit{第16页 —— 2026年1月12日（夜）}}

\noindent 交付之后我回了家。坐在床上。看着自己的手。47年。违背我意愿地活了47年。1979年我写给玛丽亚的信里说，我会保护其他人。说我不会让日记传下去。我为此相信了几十年。我相信自己掌控着局面。

\noindent 我从未掌控过。一次也没有。

\noindent 是饥饿让我活着。不是出于善意。不是出于偶然。它需要一个投递人。一个能行走在活人之间的人，有证件，有身份，看上去像个人。一个在2026年走进公共档案馆前台而不引起怀疑的人。一个81岁的老头，穿着磨旧的西装，腋下夹着一个包裹。隐形。完美。

\noindent 它从我1979年读日记那天起就计划好了。从若阿金1919年写下第62页那天起。从1853年那道壕沟开始搏动那天起。

\noindent 循环没有问过我。循环只是使用了我。

\noindent 我以为我在保护玛丽亚。以为我在斩断那条锁链。但那条锁链从来不取决于我。我只是一个环。而环不作选择。环只负责扣住。负责传递。直到轮到另一个环。

\noindent 玛丽亚死于2004年。我没有去葬礼。不是因为不想。是因为\textit{不能}。葬礼那天，我在另一座城市的一条长椅上醒来，衣服换过了，62天的记忆不见了。它把我带走了。它不会让任何人打断这次交付。

\noindent 对不起，玛丽亚。对不起，将要读到这些的医生。对不起，队伍里的下一个。

\noindent 我没有选择。

\noindent 你们也不会有。

\noindent \textit{档案员注: 这叠笔记于2026年2月在一个封口的牛皮纸信封中被发现，与患者1,260的日记同放在一个箱子里。最初几页的笔迹沉稳，是医生的字。最后几页的笔迹颤抖、不规则，其笔画令人想起——但并不完全等同于——若阿金·安东尼奥·科雷亚的笔迹。最后一则笔记的日期（2026年1月12日）与日记首次被归还米纳斯吉拉斯州公共档案馆的日期吻合。}
\noindent 原始日记的第62页——若阿金在那里写下了1853年那实体的起源——有一处水渍，部分遮住了文字。\textquotedblleft{}1853\textquotedblright{}这个日期被写了三次，笔迹各不相同——仿佛在漫长的时间里被不同的手反复确认过。最后一次确认是医生的字。

\vspace{0.5em}
\noindent\textit{\textbf{↪ 路线:} \rota{DOC-01} \rotasep \rota{LUC-02} \rotasep \rota{DOC-27}}
